\section{Implementation and Optimization}
\label{sec:implementation}

This section describes the Rust implementation used in the experiments and
the sequence of optimizations that led to the final evaluator.  The goal is
not merely to present the fastest implementation, but to separate
algorithmic structure from low-level engineering effects and to document
which optimization ideas did and did not improve performance on the tested
Apple M1 platform.

\subsection{Implementation setting}

The implementation is written in Rust and uses Rayon for CPU parallelism.
All basis implementations in a given benchmark use the same finite-field
backend.

The experiments in this paper use the Goldilocks field
\[
\mathbb{F}_p,
\qquad
p=2^{64}-2^{32}+1.
\]
Field elements are stored canonically in a single \(64\)-bit machine word.

The optimized Boolean-kernel evaluator operates on coefficient vectors of
length
\[
N=2^m.
\]
The mathematical fold is
\[
F_t(a,b)=t(a+b)+b.
\]
The implementation therefore has two distinct optimization layers:

\begin{enumerate}
    \item optimize the field-level realization of \(F_t\);
    \item organize the binary tree so that memory traffic, allocation, and
    synchronization are minimized.
\end{enumerate}

\subsection{Baseline serial implementation}

The direct serial implementation copies the coefficient vector into a work
buffer and performs the tree level by level.

If the active prefix has length \(L\), one level computes
\[
\frac{L}{2}
\]
folds
\[
\texttt{buf[j]}
\leftarrow
F_t(\texttt{buf[2j]},\texttt{buf[2j+1]}),
\]
then replaces \(L\) by \(L/2\).

This implementation closely follows the recurrence of
\cref{sec:evaluation} and is used as a correctness reference.

\subsection{Coarse-grained subtree parallelism}

A first parallel implementation may assign all independent folds at a tree
level to Rayon and synchronize before proceeding to the next level.  This
exposes maximal parallelism at the bottom of the tree, but it also creates a
global synchronization point at every level.

Instead, the optimized implementation decomposes the coefficient vector
into contiguous power-of-two blocks of size
\[
B.
\]
Each block represents a complete local subtree.  Rayon distributes those
subtrees across workers, and each worker reduces an entire subtree locally
before returning its root.

The remaining roots form a much smaller top tree, which is finished
serially.

For
\[
N=2^{20},
\qquad
B=4096=2^{12},
\]
the input is split into
\[
\frac{2^{20}}{2^{12}}=256
\]
independent local subtrees.

This organization removes the need for a global barrier after each of the
first twelve tree levels.

\subsection{Fast Goldilocks multiplication}

A generic implementation may multiply two canonical \(64\)-bit residues
into a \(128\)-bit product and reduce the result modulo
\[
p=2^{64}-2^{32}+1.
\]

The Goldilocks modulus permits a more specialized reduction.  Write
\[
x= \ell + 2^{64}h,
\]
where
\[
0\le \ell,h<2^{64}.
\]
Let
\[
\varepsilon = 2^{32}-1.
\]
Since
\[
2^{64}\equiv \varepsilon \pmod p,
\]
we have
\[
x\equiv \ell+h\varepsilon\pmod p.
\]

Split
\[
h=h_{\mathrm{lo}}+2^{32}h_{\mathrm{hi}},
\]
with
\[
0\le h_{\mathrm{lo}},h_{\mathrm{hi}}<2^{32}.
\]
Using
\[
2^{32}\varepsilon
=
2^{64}-2^{32}
\equiv -1
\pmod p,
\]
we obtain
\[
h\varepsilon
=
h_{\mathrm{lo}}\varepsilon
+
2^{32}h_{\mathrm{hi}}\varepsilon
\equiv
h_{\mathrm{lo}}\varepsilon-h_{\mathrm{hi}}
\pmod p.
\]
Therefore
\[
\boxed{
x
\equiv
\ell-h_{\mathrm{hi}}
+
h_{\mathrm{lo}}(2^{32}-1)
\pmod p.
}
\]

This form avoids the repeated high-half multiplications generated by the
more generic reduction routine.

On the tested Apple M1 system, the multiplication microbenchmark improved
from approximately
\[
2.2372\ \mathrm{ns}
\]
to
\[
2.0559\ \mathrm{ns},
\]
an improvement of roughly
\[
8\%.
\]

\subsection{Fused kernel fold}

The direct fold
\[
F_t(a,b)=t(a+b)+b
\]
could be implemented as one modular addition, one modular multiplication,
and a second modular addition.

The optimized implementation instead fuses the final \(+b\) into the
\(128\)-bit product before Goldilocks reduction.

Let
\[
s=(a+b)\bmod p.
\]
The implementation forms the full product
\[
t s
\]
as a \(128\)-bit value, adds \(b\) to that \(128\)-bit value with carry,
and only then applies the fast Goldilocks reduction.

This preserves the mathematical value
\[
t(a+b)+b
\pmod p
\]
while avoiding a second complete modular-addition sequence after the
multiplication.

The isolated fold microbenchmark improved from approximately
\[
4.0591\ \mathrm{ns}
\]
to
\[
3.3585\ \mathrm{ns},
\]
corresponding to an improvement of about
\[
17.3\%.
\]

\subsection{Bounds-check elimination}

After arithmetic optimization, the inner tree loop became short enough
that Rust bounds checks were visible in the generated assembly.

The safe indexed loop
\[
\texttt{buf[2j]},\qquad
\texttt{buf[2j+1]},\qquad
\texttt{buf[j]}
\]
generated two conditional bounds-check branches per fold iteration on the
tested compiler target.

The optimized implementation uses a small audited helper based on raw
pointers.  For an active length \(L\), the loop is executed only for
\[
0\le j<\frac{L}{2}.
\]
Therefore
\[
2j+1<L
\]
and
\[
j<\frac{L}{2}\le L.
\]
The output location of iteration \(j\) also cannot overwrite a future input
location, because for every future index \(k>j\),
\[
2k>j.
\]

Thus the in-place access pattern is valid, and the unchecked helper removes
the per-iteration bounds checks while preserving the same sequential
semantics.

For a \(4096\)-element level, the measured runtime improved from
\[
6.9671\ \mu\mathrm{s}
\]
to
\[
6.3575\ \mu\mathrm{s},
\]
an improvement of approximately
\[
8.75\%.
\]

\subsection{Direct first-level reduction}

The first subtree implementation began each local job by copying all
\(B\) input coefficients into a private mutable buffer:
\[
B
\longrightarrow
B.
\]
It then applied the first tree level:
\[
B
\longrightarrow
\frac{B}{2}.
\]

The optimized implementation eliminates this full copy.  It reads pairs
directly from the immutable input slice and writes the first-level outputs
into a newly allocated half-size buffer:
\[
\boxed{
B
\longrightarrow
\frac{B}{2}.
}
\]

For
\[
B=4096,
\]
this changes the local scratch allocation from \(4096\) field elements to
\(2048\) field elements and eliminates the initial full-block copy.

At the first-stage microbenchmark level, the measured runtime improved from
\[
6.9031\ \mu\mathrm{s}
\]
to
\[
6.1571\ \mu\mathrm{s},
\]
an improvement of approximately
\[
10.8\%.
\]

At the full-evaluator level, this optimization reduced the best measured
runtime from approximately
\[
1.0179\ \mathrm{ms}
\]
to
\[
0.93533\ \mathrm{ms}.
\]

\subsection{Worker-local scratch reuse}

Even after direct first-level reduction, creating a new half-size vector
for every subtree still incurs repeated allocation overhead.

The final implementation uses Rayon worker-local state.  Each worker owns
a reusable scratch vector and processes multiple subtrees with that same
allocation.

This differs from an earlier unsuccessful global-buffer strategy:
the source coefficient array remains immutable and all workers still read
their assigned input subtrees directly in parallel.  Only the private
scratch allocation is reused.

In the worker-scratch optimization experiment, the same-run baseline required
\[
945.77\ \mu\mathrm{s},
\]
whereas worker-local scratch reuse required
\[
892.57\ \mu\mathrm{s}.
\]
This is the best observed kernel-only optimization run; the final
cross-basis benchmark is reported separately in \cref{sec:results}.
This corresponds to an improvement of approximately
\[
5.6\%
\]
in that run.

The resulting throughput was approximately
\[
1.175\ \mathrm{Gelem/s}.
\]

\subsection{Subtree-size tuning}

The subtree size \(B\) trades scheduling overhead against local work size
and available parallelism.

After the arithmetic and memory optimizations above, we re-ran the block
size sweep because the optimum found for an earlier, slower implementation
need not remain optimal.

A broad sweep over
\[
B\in
\{1024,2048,4096,8192,16384\}
\]
identified
\[
2048,\ 4096,\ 8192
\]
as the competitive region.

A longer controlled comparison gave:

\begin{center}
\begin{tabular}{lr}
\toprule
Block size \(B\) & Median runtime \\
\midrule
2048 & \(1.0151\ \mathrm{ms}\) \\
4096 & \(920.45\ \mu\mathrm{s}\) \\
8192 & \(917.46\ \mu\mathrm{s}\) \\
\bottomrule
\end{tabular}
\end{center}

The difference between \(4096\) and \(8192\) was only about
\[
0.32\%,
\]
which is inside the observed run-to-run noise.  We therefore retain
\[
\boxed{B=4096}
\]
as the conservative default.

\subsection{Optimization progression}

The approximate progression of the best subtree-parallel implementation at
\[
N=2^{20}
\]
was
\[
1.2530\ \mathrm{ms}
\rightarrow
1.2007\ \mathrm{ms}
\rightarrow
1.1619\ \mathrm{ms}
\rightarrow
1.0179\ \mathrm{ms}
\rightarrow
0.93533\ \mathrm{ms}
\rightarrow
\boxed{0.89257\ \mathrm{ms}}.
\]

Relative to the earlier \(1.2530\)-ms subtree baseline, the final observed
runtime is lower by approximately
\[
28.8\%.
\]

The sequence is important: no single optimization accounts for the full
gain.  Arithmetic, compiler-generated checks, memory traffic, and allocation
all become visible as preceding bottlenecks are removed.

\subsection{Rejected optimization attempts}

Several plausible optimizations did not improve performance and were not
included in the final evaluator.

\paragraph{Global shared-buffer reuse.}
A version that copied the complete coefficient vector into one shared
mutable buffer and then reduced disjoint chunks avoided per-subtree
allocation, but the initial global copy serialized too much memory traffic.
It was slower than the local-subtree design.

\paragraph{Manual four-way loop unrolling.}
Explicitly computing four independent folds per source loop iteration did
not improve the generated scalar code enough to offset the additional
register pressure and code size.

\paragraph{Alternative modular addition.}
A manually derived Goldilocks addition sequence was compared with the
compiler-generated implementation.  The compiler's existing sequence was
already shorter, so the alternative was rejected before integration.

\paragraph{Three-level tree fusion.}
A local \(8\to4\to2\to1\) fused tree reduces intermediate memory traffic,
but it also increases dependency depth and register pressure.  The fused
version was slower than the ordinary level-by-level local reduction.

\paragraph{Two-level tree fusion.}
We separately tested a smaller
\[
4\to2\to1
\]
fusion.  The current two-level path required approximately
\[
9.7019\ \mu\mathrm{s},
\]
while the fused version required
\[
9.8528\ \mu\mathrm{s}.
\]
Thus the fused variant was approximately
\[
1.6\%
\]
slower in this benchmark.

These negative results suggest that on the tested Apple M1 scalar path,
reducing memory traffic is not automatically beneficial when it lengthens
the dependency chain around the already optimized field fold.

\subsection{Final implementation}

The final evaluator used for the headline benchmark combines:

\begin{enumerate}
    \item coarse-grained subtree parallelism;
    \item block size \(B=4096\);
    \item specialized Goldilocks multiplication reduction;
    \item fused kernel-fold reduction;
    \item unchecked inner-level indexing with explicit safety invariants;
    \item direct first-level reduction into a half-size scratch buffer;
    \item Rayon worker-local scratch reuse;
    \item serial completion of the small top tree.
\end{enumerate}

All of these optimizations preserve the exact fold graph derived in
\cref{sec:evaluation}.  They change only how the arithmetic is implemented
and scheduled.
